\uuid{KMX6}
\titre{Vecteurs}
\auteur{bodin ; hanani ; mzari}
\organisation{exo7}
\licence{CC BY-NC-SA 4.0}
\datecreate{2020-01-20}
\niveau{L1}
\module{Géométrie}
\chapitre{Géométrie du plan}
\sousChapitre{Distance, norme, produit scalaire}
\difficulte{1}
\theme{vecteurs}
\exo7qcmid{lille-1.geometrie-plane.007}

\contenu{
\begin{question}
\qclassification{140.01,140.02}
\qtags{motcle=vecteurs}
On considère les points $A(-1,-1)$, $B(-1,1)$, $C(1,2)$ et $D(1,0)$. Quelles sont les bonnes réponses ?
\begin{answers}
    \bad{Les droites $(AB)$ et $(CD)$ sont sécantes.}
    \bad{Les droites $(AB)$ et $(CD)$ sont perpendiculaires.}
    \good{Les droites $(AB)$ et $(CD)$ sont parallèles.}
    \good{$(ABCD)$ est un parallélogramme.}
\end{answers}
\begin{explanations}
On a $\overrightarrow{AB}=(0,2)=-\overrightarrow{CD}$, donc les droites $(AB)$ et $(CD)$ sont parallèles. De plus, $AB=CD$, donc $(ABCD)$ est un parallélogramme.
\end{explanations}
\end{question}
}
